\section{Experience}
\subsection{Core Contributor -- Unity/C\# Developer}{\href{https://www.playdreamforge.com/}{DreamForge}}{December 2025 - February 2026}
\begin{itemize}
    \item Core contributor to an AI-driven game engine that procedurally generates 3D game worlds from text prompts.
    \item Implemented a procedural \textbf{wall cutout system} using CSG boolean subtraction with four silhouette strategies (convex hull, radial ray casting, voxel tracing, alpha-shape), connected-component filtering, and ear-clipping triangulation.
    \item Built \textbf{CI/CD pipelines} using \textbf{Docker}, \textbf{Azure Container Instances}, \textbf{.NET}, \textbf{GitHub Actions}, and \textbf{Slack} webhooks for automated exception analysis, PR creation, and deployment notifications.
    \item Architected a multi-mode \textbf{camera system} using Cinemachine with adaptive framing, multi-candidate occlusion-aware selection, bone-based NPC framing, spherecast fading, and room-geometry yaw constraints.
    \item Designed a polygon-based \textbf{room shape Generator} with grid-based flood fill validation.

\end{itemize}

\subsection{Research Assistant \& Teaching Assistant}{York University, BioMotion Lab}{January 2024 - December 2025}
\begin{itemize}
    \item Conducting VR and motion capture experiments using \textbf{Meta Quest}, \textbf{ArKit}, \textbf{Unity3D}, and \textbf{Unreal Engine}, exploring depth perception by isolating stereopsis and motion parallax.
    \item \textbf{Presented research on illusory parallax in VR} at the \textbf{47th European Conference on Visual Perception (ECVP 2025)} in Mainz, Germany.
    \item Teaching Object-Oriented Programming (\textbf{Java}) for 5 semesters, emphasizing system design patterns and optimization.
\end{itemize}

\subsection{Software Developer}{TectoTrack}{August 2023 - February 2025}
\begin{itemize}
    \item Built simulation frameworks for high-traffic environments such as airports and malls, focusing on scalability \& agent realism.
    \item Designed fault-tolerant pathfinding algorithms to support dynamic and adaptive crowd movement.
    \item Developed real-time logging and monitoring dashboards to track system health and agent behavior metrics.
\end{itemize}

\subsection{Lead Unity3D Game Developer}{\href{https://techuonthechair.com}{Techu}}{January 2023 - April 2024}
\begin{itemize}
    \item \textbf{Led} the development of "Techu," a cross-platform 2D card game using \textbf{DQN} + \textbf{Monte Carlo Tree Search} for AI agents, and \textbf{PlayFab Azure} for matchmaking and data persistence.
    \item Integrated \textbf{Firebase} monitoring for real-time system diagnostics and automated crash alerts.
\end{itemize}

\subsection{Volunteer Full Stack Developer}{IAESTE}{January 2022 - December 2022}
\begin{itemize}
    \item Maintained and enhanced IAESTE’s web platform using \textbf{Vue.js} and \textbf{Django}, improving performance and user experience.
\end{itemize}

\subsection{Technical Staff Member \& Teacher Assistant}{Amirkabir University of Technology (AUT)}{September 2017 - August 2022}
\begin{itemize}
    \item Technical staff for AUT Game Development Events and TA for \textit{AI}, \textit{Advanced Java}, \textit{C Programming}, and \textit{Operating Systems}.
\end{itemize}

\subsection{Software Engineer Intern}{Sepantab (Beh Andishan Noavar)}{June 2021 - August 2021}
\begin{itemize}
    \item Developed \textit{Stone Thrower}, a local multiplayer \textbf{IoT-based WebGL} game using Unity3D and networked WebSockets.
    \item Designed resilient networking infrastructure using \textbf{C\#}, \textbf{gRPC}, \textbf{Protocol Buffers}, and \textbf{MQTT}.
\end{itemize}


% \subsection{Game Developer}{Pherma}{January 2019 – April 2020}
% \begin{itemize}
%     \item Worked as a game developer in a friendly team, developing \textbf{hyper-casual} and \textbf{adventure games} for iOS, Android and Windows platforms using \textbf{Unity} and \textbf{C\#}.
% \end{itemize}
